\begin{textsample}{POS Dim 3 – pe – Score 18.00 – pe/pe\_inf\_15.txt}  \label{ex:f3_pos_274}
2.2 PRINCÍPIOS NORTEADORES DA EDUCAÇÃO \textbf{INFANTIL} Os Princípios éticos correspondem aos valores relacionados à autonomia, à responsabilidade, à solidariedade e ao respeito ao bem comum, ao meio ambiente e às diferentes culturas, identidades e singularidades, através das propostas pedagógicas que assegurem às \textbf{crianças} a manifestação de seus interesses, \textbf{desejos} e \textbf{curiosidade} em práticas educativas que valorizem suas produções, individuais e coletivas.
As \textbf{instituições} de Educação \textbf{Infantil} devem oportunizar às \textbf{crianças} experiências que : - Ampliem a visão de mundo e de si próprias, advindas de diferentes tradições culturais. - Promovam atitudes de respeito e solidariedade, fortalecendo a autoestima e os vínculos afetivos de todas as \textbf{crianças}, combatendo preconceitos que incidem sobre as diferentes formas dos seres humanos se constituírem como pessoas. - Respeitem todas as formas de vida, o \textbf{cuidado} de seres vivos e a preservação dos recursos naturais. - Favoreçam a aquisição de valores como os da inviolabilidade da vida humana, a liberdade e a integridade individuais, a igualdade de direitos de todas as pessoas, a igualdade entre homens e mulheres, assim como a solidariedade com grupos enfraquecidos e vulneráveis política e economicamente.
Os Princípios políticos afirmam os direitos de cidadania, do exercício da criticidade e do respeito à ordem democrática.
Esses valores, na Educação \textbf{Infantil}, ganham destaque, com base na formação participativa e crítica das \textbf{crianças}, em práticas educativas que : - Possibilitem expressar sentimentos, ideias, questionamentos, comprometidos com a busca do bem-estar coletivo e individual, com a preocupação com o outro e com a coletividade. - Criem condições para que a \textbf{criança} aprenda a opinar e a considerar os sentimentos e a opinião dos outros sobre um acontecimento, uma reação afetiva, uma ideia, um conflito.
Os Princípios estéticos são referentes aos valores da sensibilidade, da criatividade, da ludicidade e da diversidade de manifestações artísticas e culturais.
Eles são evidenciados na forma de se relacionar e de ver o mundo das \textbf{crianças} e se manifestam nas práticas educativas em situações que : - Valorizem o ato criador e a construção pelas \textbf{crianças} de respostas singulares, garantindo -lhes a participação em diversificadas experiências. - Organizem um cotidiano de situações \textbf{agradáveis}, estimulantes, que desafiem o que cada \textbf{criança} e seu grupo de \textbf{crianças} já sabem sem ameaçar sua autoestima nem promover competitividade. - Ampliem as possibilidades da \textbf{criança} de \textbf{cuidar} e ser \textbf{cuidada}, de se expressar, comunicar e criar, de organizar pensamentos e ideias, de conviver, \textbf{brincar} e trabalhar em grupo, de ter iniciativa na busca de resolver problemas e conflitos que se apresentam às mais diferentes idades.
Os princípios apresentados, presentes em um trabalho pedagógico, orientam um currículo para a Educação \textbf{Infantil} que permita ao professor “ refletir, selecionar, organizar, planejar, mediar e monitorar o conjunto de práticas e interações, garantindo a pluralidade de situações que promovam o desenvolvimento pleno das \textbf{crianças} ” ( BNCC, 2017, p. 37 ).
No currículo, os Princípios de Educação \textbf{Infantil} estão materializado nos direitos de CONVIVER, \textbf{BRINCAR}, PARTICIPAR, EXPLORAR, EXPRESSAR E CONHECER -SE, esses direitos devem estar expressos nos objetivos de aprendizagem e desenvolvimento, através de experiências que sejam significativas para a \textbf{criança}.
Os direitos² de aprendizagem e desenvolvimento significam : - CONVIVER democraticamente com outras \textbf{crianças} e \textbf{adultos}, com eles se relacionar e partilhar distintas situações, utilizando diferentes linguagens, ampliando o conhecimento de si e do outro, o respeito em relação à natureza, à cultura e às diferenças entre as pessoas. - \textbf{BRINCAR} cotidianamente de diversas formas, em diferentes espaços e tempos, com diferentes \textbf{parceiros} \textbf{adultos} e \textbf{crianças}, ampliando e diversificando as culturas \textbf{infantis}, seus conhecimentos, sua imaginação, sua criatividade, suas experiências emocionais, corporais, \textbf{sensoriais}, expressivas, cognitivas, sociais e \textbf{relacionais}. - PARTICIPAR ativamente, junto aos \textbf{adultos} e às outras \textbf{crianças}, tanto do planejamento da gestão da escola quanto da realização das atividades da vida cotidiana, da escolha das \textbf{brincadeiras}, dos materiais e dos ambientes, desenvolvendo linguagens e elaborando conhecimentos, decidindo e se posicionando. - EXPLORAR movimentos, \textbf{gestos}, \textbf{sons}, palavras, emoções, transformações, relacionamentos, histórias, objetos, elementos da natureza - no contexto urbano e do campo -, espaços e tempos das \textbf{instituições}, interagindo com diferentes grupos e ampliando seus saberes, linguagens e conhecimentos. - EXPRESSAR, como sujeito criativo e sensível, com diferentes linguagens, sensações corporais, necessidades, opiniões, sentimentos e \textbf{desejos}, pedidos de \textbf{ajuda}, narrativas, \textbf{registros} de conhecimentos elaborados a partir de diferentes experiências, envolvendo tanto a produção de linguagens quanto a fruição das artes em todas as suas manifestações. - CONHECER -SE e construir sua identidade pessoal, social e cultural, constituindo uma imagem positiva de si e de seus grupos de pertencimento, nas diversas experiências de \textbf{cuidados}, interações e \textbf{brincadeiras} vivenciadas na \textbf{instituição} de Educação \textbf{Infantil}. ² Fonte : BNCC, 2ª versão preliminar, 2016, p. 63.
Compreendendo que esses direitos apresentam uma visão integralizada do desenvolvimento \textbf{infantil}, este currículo, em consonância com a BNCC, é norteador para elaborações e adequações das propostas pedagógicas de redes municipais e de \textbf{instituições} de Educação \textbf{Infantil} de Pernambuco, como também de programas e projetos pedagógicos para essa etapa de ensino.
Assim, o currículo na Educação \textbf{Infantil} prioriza a formação identitária, a ludicidade, a autonomia, a autoestima, a cooperação, as interações e as \textbf{brincadeiras} no processo de aprendizagem e desenvolvimento da \textbf{criança}, sendo o \textbf{brincar} aspecto significativo de possibilidades na criação de situações cotidianas que permitam a construção da sua identidade, da imagem de si mesmo e do mundo em que vive.

% matched lemmas: adulto, agradável, ajuda, brincadeira, brincar, criança, cuidado, cuidar, curiosidade, desejo, gesto, infantil, instituição, parceiro, registro, relacional, sensorial, som
\end{textsample}
